% ID: SEL-D0069-OB-E03
% ESTADO: validada. S069.
\textbf{Parte I.} Sean $E$ «estudia Estadística», $H$ «es hombre» y $M$ «es mujer».

\textbf{a)} Por probabilidad total,
\[
P(H)=0{,}70\cdot0{,}40+0{,}30\cdot0{,}70=\boxed{0{,}49}.
\]
\textbf{b)} Como $P(M)=0{,}51$ y $P(E\cap M)=0{,}70\cdot0{,}60=0{,}42$, por Bayes
$P(E\mid M)=0{,}42/0{,}51=\boxed{14/17}$.

\textbf{Parte II. a)} $\hat p=75/300=0{,}25$. Los recuentos $75$ y $225$ permiten la aproximación normal. Al $94\%$ queda $0{,}03$ en cada cola, con $z_{0{,}97}\approx1{,}88079$:
\[
0{,}25\pm 1{,}88079\sqrt{\frac{0{,}25\cdot0{,}75}{300}}
\approx\boxed{[0{,}2030,\ 0{,}2970]}.
\]
\textbf{b)} El error máximo es la semiamplitud, $\boxed{0{,}0470}$ aproximadamente, o $4{,}70$ puntos porcentuales.
